\documentclass[11pt]{article}
\textwidth 15.9cm
\textheight 23. cm
\addtolength{\oddsidemargin}{-11.5mm}
\addtolength{\evensidemargin}{-11.5mm}
\addtolength{\topmargin}{-19.0mm}

%\usepackage{amsmath}
\usepackage[ansinew]{inputenc}
\usepackage[bbgreekl]{mathbbol}
\usepackage{amsmath,amssymb}
\usepackage{graphicx}
\usepackage{hyperref}
\usepackage{url}
\usepackage[numbers,square]{natbib}
\usepackage{multirow}
\usepackage{bbm}
\usepackage{authblk}
\usepackage{subcaption}
\renewcommand{\Affilfont}{\footnotesize}
%\renewcommand{\Authfont}{\normalsize}

\usepackage{stmaryrd}
\usepackage{comment}
\usepackage{showlabels}
\usepackage{soul}
%\usepackage[dvipsnames]{xcolor}

%%%%%%%% MACROS
\include{gempic_spectral_macros}

\newcommand{\hi}{\hat{i}}
\newcommand{\hj}{\hat{j}}
\newcommand{\hk}{\hat{k}}

% \newcommand{\arr}[1]{{\bs{\mathsf{#1}}}}

% matrices
\newcommand{\hodgeZero}{\mathbb{H}_0}
\newcommand{\hodgeOne}{\mathbb{H}_1}
\newcommand{\hodgeTwo}{\mathbb{H}_2}
\newcommand{\hodgeThree}{\mathbb{H}_3}
\newcommand{\hodgeDualZero}{\tilde{\mathbb{H}}_0}
\newcommand{\hodgeDualOne}{\tilde{\mathbb{H}}_1}
\newcommand{\hodgeDualTwo}{\tilde{\mathbb{H}}_2}
\newcommand{\hodgeDualThree}{\tilde{\mathbb{H}}_3}
\newcommand{\hodgeDualZeroB}{\mathbb{H}^*_0}
\newcommand{\hodgeDualOneB}{\mathbb{H}^*_1}
\newcommand{\hodgeDualTwoB}{\mathbb{H}^*_2}
\newcommand{\hodgeDualThreeB}{\mathbb{H}^*_3}

\newtheorem{theorem}{Theorem}
\newtheorem{lemma}{Lemma}
\newtheorem{remark}{Remark}
\newtheorem{definition}{Definition}
\newtheorem{assumption}{Assumption}

\begin{document}

\section{Time loops}

\subsection{Hamiltonian splitting for Vlasov-Maxwell}

The semi-discrete Vlasov-Maxwell equations on staggered grids form a noncanonical Hamiltonian sytem of the form $\fract{\arrU}{t}=\mathcal{J}(\arrU)\nabla\cH$, where $\arrU$ stacks all
the particle positions $\arrX$, velocities $\arrV$ as well as the electric field $\arrD$ and magnetic field $\arrB$.
The Hamiltonian is 
\begin{equation}
    \cH =   \tfrac 12 \arrV^\top \matW_m \arrV + \tfrac 12 \arrD^\top  \hodgeDualTwo \arrD + \tfrac 12 \arrB^\top \hodgeTwo \arrB,
\end{equation}
where $\matW_m$ is a diagonal matrix with the mass of the corresponding particle on the diagonal. We also denote by $\arrE =  \hodgeDualTwo \arrD$ and $\tilde{\arrH}=\hodgeTwo \arrB$.
The hamiltonian consists of three parts, the kinetic energy $\cH_p = \tfrac 12 \arrV^\top \matW_m \arrV$, the electric energy $\cH_e = \tfrac 12 \arrD^\top  \hodgeDualTwo \arrD$
and the magnetic energy $\cH_b= \tfrac 12 \arrB^\top \hodgeTwo \arrB$.

We solve the equations using a hamiltonian splitting procedure, which yields three subsystems. We denote by $\bE$, $\bB$ the smoothed electric and magnetic fields used for particle mesh interactions.
\begin{enumerate}
    \item Keeping only $\cH_b$:
    $$
    \begin{aligned}
        \fract{\bx_p}{t} &= 0, \\
        \fract{\bv_p }{t} &= 0 \\
 	    \fract{\tilde{\arrD}}{t} &=  \tilde{\mathbb{C}}\tilde{\arrH}, \\
 	    \fract{\arrB }{t}&= 0.
    \end{aligned}
    $$
    Since $\arrB$ (and hence $\tilde{\arrH}$) are constant in this step, the exact solution of the system is straightforward.
    \item Keeping only $\cH_e$:
    $$
    \begin{aligned}
        \fract{\bx_p}{t} &= 0, \\
        \fract{\bv_p }{t} &= \frac {q_p}{m_p} \bE(\bx_p) \\
 	    \fract{\arrD}{t} &=  0, \\
 	    \fract{\arrB }{t}&= - \mathbb{C}\arrE.
    \end{aligned}
    $$
    Since $\bx_p$ and $\arrD$ (and hence $\arrE$) are constant in this step, the exact solution of the system is straightforward.
    \item Keeping only $\cH_p$:
    $$
    \begin{aligned}
        \fract{\bx_p}{t} &= \bv_p, \\
        \fract{\bv_p }{t} &= \frac {q_p}{m_p} \bv_p\times\bB(\bx_p) \\
 	    \fract{\tilde{\arrD}}{t} &=  -\tilde{\arrJ}, \\
 	    \fract{\arrB }{t}&= 0.
    \end{aligned}
    $$
    where $\tilde{\arrJ}$ depends on all the particles positions and velocities.
    This system cannot be solved explicitly in general, but we can further split the kinetic energy into its three components
    $\cH_p = \cH_{p,x} + \cH_{p,y} + \cH_{p,z}$. We write $\bx_p=(x_p,y_p,z_p)$, $\bv_p=(v_{p,x},v_{p,y},v_{p,z})$ 
    \begin{enumerate}
        \item Keeping only $\cH_{p,x}$, and writing only the terms with non vanishing right-hand side
        $$
    \begin{aligned}
        \fract{x}{t} &= v_{p,x}, \\
        \fract{v_{p,y} }{t} &= - \frac {q_p}{m_p} v_{p,x} B_z(x(t),y,z) \\
        \fract{v_{p,z} }{t} &=  \frac {q_p}{m_p} v_{p,x} B_y(x(t),y,z) \\
 	    \fract{\tilde{\arrD}_x}{t} &=  -\tilde{\arrJ}_x.
    \end{aligned}
    $$
    Using that $v_{p,x}=\fract{x}{t}$, this can be integrated exactly by using the primitive in $x$ of $\bB$ and $\bJ_x$
    \item Keeping only $\cH_{p,y}$, and writing only the terms with non vanishing right-hand side
        $$
    \begin{aligned}
        \fract{y}{t} &= v_{p,y}, \\
        \fract{v_{p,x} }{t} &= \frac {q_p}{m_p} v_{p,y} B_z(x,y(t),z) \\
        \fract{v_{p,z} }{t} &= - \frac {q_p}{m_p} v_{p,y} B_x(x,y(t),z) \\
 	    \fract{\tilde{\arrD}_y}{t} &=  -\tilde{\arrJ}_y.
    \end{aligned}
    $$
    Using that $v_{p,y}=\fract{y}{t}$, this can be integrated exactly by using the primitive in $y$ of $\bB$ and $\bJ_y$.
    \item Keeping only $\cH_{p,z}$, and writing only the terms with non vanishing right-hand side
    $$
\begin{aligned}
    \fract{z}{t} &= v_{p,z}, \\
    \fract{v_{p,x} }{t} &= - \frac {q_p}{m_p} v_{p,z} B_y(x,y,z(t)) \\
    \fract{v_{p,y} }{t} &=  \frac {q_p}{m_p} v_{p,z} B_y(x,y,z(t)) \\
     \fract{\tilde{\arrD}_z}{t} &=  -\tilde{\arrJ}_z.
\end{aligned}
$$
Using that $v_{p,z}=\fract{z}{t}$, this can be integrated exactly by using the primitive in $z$ of $\bB$ and $\bJ_z$.
    \end{enumerate}

\end{enumerate}

\textbf{Remark:} In two space dimensions keeping all velocity components and all components of the fields, $v_{p,z}$ is constant in this step and $\bB$ does not depend on $z$, so that 
all equations can be integrated exactly in a straightforward manner without using the primitive,
including the last equation using the definition of $\tilde{\arrJ}_z$.

In one space dimension, the same also applies to the previous step.

\end{document}